% from ega4-18-fr.tex
\oldpage[IV-4]{154}

est un corps et que $f$ est plat, $B$ est régulier (0, 17.3.3); mais en outre (6.1.1.1), on a $\dim(B)=\dim(A)+\dim(B\otimes_A k)=\dim(A)=1$, donc $B$ est aussi un anneau de valuation discrète. Mais alors $X-\{b\}$ est réduit à un seul point, d'où le lemme (18.9.7.3) dans ce cas.

On observera que cela prouve l'énoncé du théorème (18.9.7) lorsqu'on suppose en outre que, dans cet énoncé, l'anneau $A$ est \emph{japonais}, car en vertu de la remarque (18.9.7.1), on est ramené au cas où $A$ est \emph{intégralement clos}, et alors, dans la réduction qui précède (18.9.7.3), l'anneau $A_1=\sh{O}_{Y,y}$ est lui aussi intégralement clos et on peut appliquer le résultat qui vient d'être démontré.

B) \emph{Cas où $\dim(A)\geq 2$.} Le résultat de (18.9.7.3) résultera alors des deux lemmes suivants:

\begin{lemma}[18.9.7.5]
\label{IV.18.9.7.5}
Soient $A$ un anneau semi-local noethérien et réduit, $A'$ sa fermeture intégrale dans son anneau total des fractions, $X=\Spec(A)$, $X'=\Spec(A')$; soient $S$ l'ensemble des points fermés de $X$ et $U=X-S$. Supposons que pour tout point $x'\in X'$ au-dessus d'un point de $S$, on ait $\dim(\sh{O}_{X',x'})\geq 2$; alors l'anneau $A^{(1)}=\Gamma(U,\sh{O}_X)$ est entier sur $A$ et est une $A$-algèbre isomorphe à une sous-algèbre de $A'$.
\end{lemma}

Rappelons que si les $\mathfrak p_i$ ($1\leq i\leq m$) sont les idéaux minimaux de $A$, $A'$ est le composé direct des clôtures intégrales $A'_i$ des anneaux intègres $A_i=A/\mathfrak p_i$, de sorte que $X'$ est la somme des schémas $X'_i=\Spec(A'_i)$. Si $X_0$ est la somme des schémas $X_i=\Spec(A_i)$ ($1\leq i\leq m$) et $U_0$ l'image réciproque de $U$ dans $X_0$, l'homomorphisme canonique $\Gamma(U,\sh{O}_X)\to\Gamma(U_0,\sh{O}_{X_0})$ est injectif puisque $X$ est réduit et le morphisme $X_0\to X$ (donc aussi le morphisme $U_0\to U$) surjectif. Comme $\Gamma(U_0,\sh{O}_{X_0})$ est composé direct des $\Gamma(U_i,\sh{O}_{X_i})$ ($1\leq i\leq m$), où $U_i$ est l'image réciproque de $U$ dans $X_i$, tout revient à prouver que pour chaque $i$, $\Gamma(U_i,\sh{O}_{X_i})$ est isomorphe à une sous-$A_i$-algèbre de $A'_i$. Autrement dit, on est ramené au cas où $A$ est un anneau semi-local noethérien intègre. Comme $X$ est réduit et le morphisme $X'\to X$ surjectif, l'homomorphisme $\Gamma(U,\sh{O}_X)\to\Gamma(U',\sh{O}_{X'})$ (où $U'$ est l'image réciproque de $U$ dans $X'$) est injectif, et tout revient à voir que l'homomorphisme canonique $A'=\Gamma(X',\sh{O}_{X'})\to\Gamma(U',\sh{O}_{X'})$ est \emph{bijectif}. Or (I, 8.2.1.1), $\Gamma(U',\sh{O}_{X'})$ est l'intersection des anneaux locaux $A'_{\mathfrak p'}$, où $\mathfrak p'$ parcourt $U'$, et en vertu de l'hypothèse, parmi ces anneaux locaux figurent tous ceux pour lesquels $\mathfrak p'$ est de hauteur 1. Mais le raisonnement de (0, 23.2.7) s'applique aussi à un anneau semi-local intègre noethérien, donc $A'$ est un \emph{anneau de Krull semi-local}, et est donc l'intersection des anneaux locaux $A'_{\mathfrak p'}$, où $\mathfrak p'$ parcourt l'ensemble des idéaux premiers de hauteur 1 de $A'$ (Bourbaki, \emph{Alg. comm.}, chap.~VII, \textsection~1, n\textsuperscript{o}~6, th.~4); \emph{a fortiori}, $A'$ est l'intersection des $A'_{\mathfrak p'}$ pour $\mathfrak p'\in U'$, ce qui achève la démonstration du lemme.

\begin{remark}[18.9.7.6]
\label{IV.18.9.7.6}
On sait (0, 23.2.5) qu'il existe une sous-$A$-algèbre finie $A''$ du corps des fractions $K$ de $A$ telle que le morphisme $\Spec(A')\to\Spec(A'')$ soit \emph{radiciel}; comme ce morphisme est aussi entier et dominant, donc surjectif (II, 6.1.10), c'est un \emph{homéomorphisme} (2.4.5). En vertu de l'interprétation géométrique de $\dim(\sh{O}_{X',x'})$ (5.1.2), on a donc, pour tout point $x''$ de $X''=\Spec(A'')$ au-dessus d'un point $x\in X$, $\dim(\sh{O}_{X'',x''})=\dim(\sh{O}_{X',x'})\geq 2$ si $x'$ est l'unique point de $X'$ au-dessus de $x$. Cela

